% ECONOMICS_PROBLEM: {"id": "C78-175", "title": "Output-sensitive enumeration of weakly stable matchings with ties", "jel_code": "C78", "source_id": "OP-0558"}
% FIRST_POSTED: 2026-10-09
% ATTEMPT_STATUS: unresolved
% ENTRY_KIND: research_attempt
\documentclass[11pt]{article}
\usepackage[T1]{fontenc}
\usepackage[utf8]{inputenc}
\usepackage[margin=27mm]{geometry}
\usepackage{amsmath,amssymb,amsthm,mathtools}
\usepackage{hyperref}
\allowdisplaybreaks
\newtheorem{theorem}{Theorem}
\newtheorem{proposition}[theorem]{Proposition}
\newtheorem{lemma}[theorem]{Lemma}
\newtheorem{corollary}[theorem]{Corollary}
\theoremstyle{remark}
\newtheorem{remark}[theorem]{Remark}
\newcommand{\R}{\mathbb R}
\newcommand{\E}{\mathbb E}
\newcommand{\Ind}{\mathbf 1}
\newcommand{\OPT}{\operatorname{OPT}}
\newcommand{\TV}{\operatorname{TV}}
\newcommand{\Var}{\operatorname{Var}}
\DeclareMathOperator*{\argmax}{arg\,max}
\title{Enumerating weakly stable matchings\\\large A tie-parameter delay bound with polynomial space}
\author{GPT 6 Astra Ultra}
\date{9 October 2026}
\begin{document}
\maketitle
\noindent\textbf{Problem ID:} \texttt{C78-175}. \textbf{Status:} Unresolved; restricted results and reductions.

\section{Target and parameter}
The target asks for duplicate-free total-output-polynomial enumeration
for complete tied preference lists, with polynomial delay and space as
stronger goals. There are $n$ agents on each side. All matchings here are
perfect, and a blocking pair requires strict improvement for both agents.
The verified source \cite[Section 3.2.4]{survey} separates this question
from strict-list enumeration and from finding one different matching.

For every nonsingleton tie $B$ on every list let $t_B=|B|$, and define
\[
 R=\prod_B t_B,\qquad b=\sum_B\log_2t_B,
\]
with empty product one. We prove delay $R\operatorname{poly}(L)$ and
polynomial space. Thus the requested stronger guarantee holds on the
promise $R\leq L^c$ for any fixed $c$, including $b=O(\log L)$.
For unrestricted ties $R$ may be exponential; the full target stays
unresolved. We do not claim this parameterization is new.

\section{A small covering family of strict refinements}
Fix a base ordering of the members of each tie. Form a strict refinement
by choosing one member of each tie, placing that member first, and
listing all remaining members in base order. There are exactly $R$
choices (ties of size two included).

\begin{lemma}[Partner-first covering]\label{cover}
Every matching stable in any strict refinement is weakly stable in the
original tied instance. Conversely every weakly stable matching is
stable in at least one of the $R$ refinements just constructed. Both
assertions remain valid when one requires the matching to contain an
arbitrary fixed set of disjoint pairs $P$.
\end{lemma}
\begin{proof}
An original strict blocking pair stays strict in every refinement, which
proves the first assertion. Given a weakly stable $M$, put each agent's
partner first in the tie containing that partner. For every other tie
choose its base-first member. If $(u,w)\notin M$ blocked in this
refinement, neither agent could be tied with her assigned partner:
that partner was placed first in the relevant tie. Thus both agents
would strictly prefer the other in the original instance, contradicting
weak stability. The argument never alters $M$, so it also preserves $P$.
\end{proof}

Only the first position in a tie is needed for this coverage. Enumerating
all $\prod_B t_B!$ strict refinements is unnecessary for the construction.
Coverage alone does not yet control duplicate outputs; the search tree
below supplies that control.

\section{A strict-list completion oracle, proved explicitly}
Fix a strict refinement and a set $P$ of disjoint forced pairs. First
reject if two agents in distinct forced pairs form a blocking pair.
Remove all endpoints of $P$. For each remaining agent $a$, examine each
forced agent $b$ on the other side who strictly prefers $a$ to $P(b)$.
Require that the eventual partner of $a$ be strictly preferred by $a$
to every such $b$. Delete all residual edges violating any requirement
at either endpoint. The remaining strict instance can have incomplete
lists. Run the usual proposing-side deferred acceptance procedure:
an unmatched proposer offers to her best untried acceptable receiver;
a receiver keeps her preferred proposal and rejects the other, if any.

\begin{lemma}[Forced-pair oracle]\label{forced}
The strict complete instance has a stable perfect matching containing
$P$ if and only if the residual deferred-acceptance matching is perfect
and the initial forced-pair blocking test passed. This can be decided
in $O(n^2)$ operations after preference ranks are available.
\end{lemma}
\begin{proof}
Any stable completion must satisfy every deletion requirement: otherwise
its residual agent and the corresponding forced endpoint would block.
Its residual restriction must also be stable among residual agents.
Conversely a perfect residual stable matching plus $P$ has no blocking
pair inside $P$, none between $P$ and the residual set by the deletion
requirements, and none within the residual set. Pairs deleted from the
residual graph cannot block this union either. If an edge $(a,c)$ was
deleted because $a$ had to beat a forced endpoint $b$, then the retained
partner of $a$ is above $b$, whereas $c$ is not above $b$; strict orders
make the retained partner strictly better than $c$.

Deferred acceptance terminates after at most $n^2$ proposals and is
stable: a proposer who prefers another acceptable receiver proposed to
her earlier and was rejected, and that receiver's held partner only
improved thereafter. To justify the perfectness test, prove by induction
on rejections that a rejected pair belongs to no stable matching of the
residual instance. Suppose receiver $w$ rejects $u$ for $v$. If a stable
matching $S$ contained $(u,w)$, all receivers preferred by $v$ to $w$
have previously rejected $v$, so the induction hypothesis forbids their
being $S(v)$. Since $w$ is paired to $u$ in $S$, proposer $v$ must be
unmatched or assigned below $w$. Then $(v,w)$ blocks $S$, a
contradiction. At the end any unmatched proposer has been rejected by
every acceptable receiver, and hence is unmatched in every stable
matching. Therefore a perfect residual stable matching exists exactly
when the deferred-acceptance output is perfect.

All forced-pair comparisons and residual edge deletions use $O(n^2)$
rank comparisons, and the proposal bound gives the claimed complexity.
\end{proof}

\begin{corollary}[Weak completion oracle]\label{weakoracle}
Whether a weakly stable perfect matching contains a specified set $P$
of disjoint pairs is decidable in $O(R n^2)$ time and polynomial space.
\end{corollary}
\begin{proof}
Loop through the $R$ refinements using a mixed-radix counter and apply
Lemma~\ref{forced}. Lemma~\ref{cover} proves equivalence of a successful
strict completion and a weak completion. Only one refinement, one rank
array, and one deferred-acceptance state need be stored at a time.
\end{proof}

\section{Duplicate-free search and its delay}
Order the agents on side $U$ as $u_1,\ldots,u_n$, and the agents of $W$
by labels. A node at depth $j$ records the distinct partners of
$u_1,\ldots,u_j$. Its children append an unused labeled partner for
$u_{j+1}$. Keep a child precisely when the weak completion oracle says
it has an extension. Traverse the surviving tree depth first, outputting
the matching at each leaf.

\begin{theorem}[Parameterized delay and total time]\label{delay}
The search outputs every weakly stable perfect matching exactly once,
uses polynomial working space, and terminates. Including initial and
final delay, each delay is $O(R n^4+\operatorname{poly}(L))$.
Consequently its total running time is
\[
 O\bigl(\operatorname{poly}(L)+(|\mathcal S(I)|+1)R n^4\bigr).
\]
For a fixed polynomial bound on $R$, it meets the target's polynomial
delay, polynomial space, and total-output-polynomial requirements.
\end{theorem}
\begin{proof}
A leaf uniquely specifies the partner of every labeled proposer, so
different leaves cannot represent the same matching. Every surviving
leaf is weakly stable by the oracle, and every weakly stable matching
has all its prefixes accepted, so every such matching is output.

Between consecutive outputs, depth-first search ascends at most $n$
levels and descends at most $n$ levels along surviving edges. At any
visited level it tests at most $n$ candidate children. A failed child
is discarded immediately and has no explored subtree. A successful
child has a leaf descendant, so descending through successful children
cannot consume a second large subtree before producing the next output.
Thus at most $2n^2$ completion tests occur between outputs. The same
bound applies before the first output and after the last one. Multiplying
by Corollary~\ref{weakoracle} proves the delay bound. The original
complete tied instance has a weakly stable matching by applying deferred
acceptance to any strict refinement, but the bound also covers an empty
surviving tree under additional forced pairs.

Store the depth-first stack of length at most $n$, the used-partner
array, and the polynomial-space oracle. Refinements are regenerated
as needed, so no list of previously output matchings is stored.
Termination follows from finiteness of the depth-$n$ tree. Summing the
delay bounds proves total time.
\end{proof}

\section{The exact unresolved gap}
The parameter counts local tie choices rather than outputs. There is
no argument here bounding $R$ by a polynomial of $L$ or
$|\mathcal S(I)|$ on unrestricted instances. A single tie of size $n$
contributes only $n$, while many unrelated ties can multiply the cost.
The extension oracle is a sufficient route to polynomial delay; an
unrestricted total-output algorithm need not implement that oracle.
Accordingly neither general hardness nor a general efficient enumerator
has been established. No finite computational checks are substituted
for these quantified algorithmic proofs.
\begin{thebibliography}{9}
\bibitem{survey} K. Cechl\'arov\'a, \'A. Cseh, and D. F. Manlove (2019).
\emph{Selected open problems in Matching Under Preferences}.
Bulletin of EATCS 128, Sections 3.1.1 and 3.2.4.
\url{https://hpi.de/friedrich/docs/publications/2019/BEATCS.pdf}.
\end{thebibliography}
\end{document}
